\uuid{Dl4F}
\exo7id{4547}
\titre{exo7 4547}
\auteur{quercia}
\organisation{exo7}
\datecreate{2010-03-14}
\isIndication{false}
\isCorrection{false}
\chapitre{Suite et série de fonctions}
\sousChapitre{Autre}
\cpge{mpsi-mp2i.fonctions-deux-variables.b, pcsi.fonctions-deux-variables.b, ptsi.fonctions-deux-variables.b, bcpst1.fonctions-deux-variables.c, mp-mpi.suites-series-fonctions.a, pc.suites-series-fonctions.a, psi.suites-series-fonctions.a, mp-mpi.series-entieres.d, pc.series-entieres.c, psi.series-entieres.c, pt.series-entieres.c}
\module{Analyse}
\niveau{L2}
\difficulte{}

\contenu{
\texte{
Pour~$y\in\R$ et~$n\in\N^*$, on pose $a_n(y) = \frac{\cos(ny)}{\sqrt n}$.
}
\begin{enumerate}
    \item \question{Déterminer le rayon de convergence de la série entière $\sum a_n(y)x^n$.}
    \item \question{Soit~$D=\{(x,y)\in\R^2,\ |x|<1\}$ et $F(x,y) = \sum_{n=1}^{+\infty} a_n(y)x^n$.
    Montrer que $F$, $\frac{\partial F}{\partial x}$ et $\frac{\partial F}{\partial y}$ existent en tout point de~$D$.}
\end{enumerate}
}
